\chapter{Confiabilidade, Observabilidade e Operacao}
\label{chap:confiabilidade_observabilidade}

\section{Escopo do capitulo}
Indice de confiabilidade, logs, traces, monitoracao e resiliencia.

\section{Objetivos de aprendizagem}
\begin{itemize}[leftmargin=*]
  \item Entender os conceitos nucleares do tema.
  \item Aplicar o tema em uma implementacao concreta do projeto.
  \item Definir evidencias objetivas de qualidade e funcionamento.
\end{itemize}

\section{Implementacao orientada por passos}
\begin{enumerate}[leftmargin=*]
  \item Preparacao do ambiente e insumos.
  \item Implementacao guiada do fluxo principal.
  \item Validacao dos resultados esperados.
  \item Registro de decisoes tecnicas e trade-offs.
\end{enumerate}

\section{Plano de testes do capitulo}
\subsection{Teste funcional}
\begin{itemize}[leftmargin=*]
  \item Definir casos de entrada e saida esperada.
  \item Verificar comportamento nominal e casos de erro.
\end{itemize}

\subsection{Teste de desempenho}
\begin{itemize}[leftmargin=*]
  \item Medir latencia ponta a ponta.
  \item Medir uso estimado de tokens por execucao.
\end{itemize}

\subsection{Teste de qualidade de resposta}
\begin{itemize}[leftmargin=*]
  \item Medir aderencia ao contexto recuperado.
  \item Medir consistencia da resposta com o objetivo do usuario.
\end{itemize}

\section{Criterios de aceite}
\begin{itemize}[leftmargin=*]
  \item Implementacao executa sem erros bloqueantes.
  \item Testes definidos no capitulo atingem resultado esperado.
  \item Evidencias de execucao sao anexadas pelos autores.
\end{itemize}

\section{Espaco para expansao iterativa}
\begin{itemize}[leftmargin=*]
  \item Notas de versao do capitulo.
  \item Melhorias planejadas.
  \item Lacunas para investigacao futura.
\end{itemize}
